\section{Scientific provenance and reproducibility}
\label{app:provenance}

The manuscript is derived from banked finite-dimensional theorem, review, and reconciliation artifacts.  The development repository is private; repository visibility is not changed by this paper-authoring lane.  The public archival release currently associated with the broader programme has Figshare DOI
\begin{center}
  \url{https://doi.org/10.6084/m9.figshare.33716677}.
\end{center}
That release is contextual only and is not represented as containing every artifact used in this manuscript.

\subsection{Load-bearing repository authorities}

\begin{table}[htbp]
\centering
\small
\caption{Load-bearing repository authorities for the present draft.  Review-only pull requests are cited as provenance and were not merged to create scientific authority.}
\label{tab:authority-provenance}
\begin{tabular}{@{}>{\raggedright\arraybackslash}p{0.19\linewidth}>{\raggedright\arraybackslash}p{0.29\linewidth}>{\raggedright\arraybackslash}p{0.44\linewidth}@{}}
\toprule
Authority & Frozen identifier & Role \\
\midrule
Microscopic predecessor & \commitsha{cb715a13254c4ff9cdc4594767e1f5cf5f52d954} & Frozen finite-model provenance used by the independent one-spin review. \\
\addlinespace
One-spin author theorem & \commitsha{76dde863ab0433907a527229acb04e23b8c5c396} & Exact one-known-matter-spin recoverability theorem inside the finite record study. \\
\addlinespace
PR \#195 hostile review & \commitsha{2f328e7d9d72a454467ddbd27ebd2fb53e5d66c4} & Independent proof-first authority for the projected-Pauli identity, Knill--Laflamme recovery, virtual-qubit factorization, sector invariance, and reconstructive closure classification. \\
\addlinespace
PR \#196 author study & \commitsha{0bbe6d7d6fceb35c429e38758eaeff1231bc06a3} & Finite interaction-closed record constructions $J_2/J_4/J_6$ at $N=8,10,12$. \\
\addlinespace
PR \#198 hostile review & scientific terminal \commitsha{a1904d70b46df79d903134089e3263742e2d0954}; handoff \commitsha{becbc82e77b5c4fe8c595ae7d77008b1bd202582} & Independent acceptance of exact finite-$N$ reconstructive autonomy with scope hardening; records the nonuniqueness of off-image CPTP extensions and rejects minimal-support extrapolation. \\
\addlinespace
PR \#199 reconciliation & \commitsha{bb17394b8ae71f547db29801e0f56c36042d3ad5} & Reconciles the dynamical-record frontier: exact reconstructive closure accepted, genuinely lossy autonomous closure and minimum-support scaling unestablished. \\
\bottomrule
\end{tabular}
\end{table}

A source-tree-level authority ledger is maintained alongside the manuscript in \texttt{AUTHORITIES.md}.  See \texttt{CLAIMS.md} for the frozen claim ceiling.  See \texttt{LITERATURE-NOTES.md} for literature status.

\subsection{Analytic versus numerical status}

The central theorem does not depend on numerical tolerance.  The identities
\begin{equation}
  P_0\sigma_i^\alpha P_0=0,
  \qquad
  P_0\ket{s}\!\bra{t}_iP_0=\frac{\delta_{st}}2P_0
\end{equation}
are proved analytically from collective-$\SU(2)$ representation theory, and the recovery map in \cref{eq:explicit-recovery} is explicit.  Likewise, singlet-sector preservation follows analytically from the scalar character of the pairwise matter couplings.  Previously banked finite-dimensional numerical checks are corroborative rather than load-bearing for these claims.

The finite record labels and support counts in \cref{tab:finite-records} are imported from the reviewed record-search artifacts.  This paper does not rerun the search, open new parameter points, optimize a recovery, or create new numerical evidence.

\subsection{Reproduction of the manuscript}

The source tree is self-contained at the LaTeX level:
\begin{verbatim}
publications/papers/known-spin-erasure/
  main.tex
  references.bib
  sections/
  appendices/
  AUTHORITIES.md
  LITERATURE-NOTES.md
  CLAIMS.md
  README.md
\end{verbatim}
A conventional BibTeX build is sufficient.  No generated figure is required for the present draft; the key recovery and closure structures are given by equations \eqref{eq:virtual-schmidt}, \eqref{eq:erasure-W-form}, and \eqref{eq:reconstructive-diagram}.  Future figures, if added, should clarify these structures rather than introduce independent scientific claims.
